% BOOK-FILE-ID: f-edcd2730cbdd
% 付録：高校範囲の解法と大学数学の見方を結ぶ5つの題材

\section{大学数学とのつながり}\label{節：大学数学とのつながり}

本編では,高校数学の範囲でたくさんの漸化式を解いてきた.
その解き方の多くは,大学で学ぶ数学から見直すと,もっと見通しよく理解できる.
この節では5つの題材を取り上げ,まず高校範囲の解法を確認してから,大学数学の見方を紹介する.
それぞれの題材で,大学数学の見方が次のどれを説明しているのかを意識して読んでほしい.
\begin{enumerate}
  \item その置き換えを\textbf{思いつく理由}
  \item 一見別々の解法に共通する\textbf{仕組み}
  \item 扱える問題の\textbf{広がり}
\end{enumerate}
高校範囲を超える内容も含むが,必要な定義はその場で述べる.
一つ一つ丁寧に追えば読めるはずである.

\subsection{基本対称式型と双曲線関数}\label{節：基本対称式型と双曲線関数}

\paragraphbf{高校範囲の解法}
\sectionref*{節：対称式を使う}では,
\[
  a_1=3,\qquad a_{n+1}=a_n^2-2
\]
を,\,$a_n=u_n+\frac{1}{u_n}$と置いて解いた.
$\left(u+\frac1u\right)^2-2=u^2+\frac{1}{u^2}$だから,\,$u_n$は1回進むごとに2乗され,
\[
  a_n=\alpha^{2^{n-1}}+\alpha^{-2^{n-1}},\qquad \alpha=\frac{3+\sqrt5}{2}
\]
を得たのだった.
しかし,なぜ「積が1になる2数の和」を思いつけるのだろうか.

\paragraphbf{双曲線関数で見る}
ここで,次の関数を導入しよう.
\[
  \cosh x=\frac{e^x+e^{-x}}{2}
\]
これを双曲線余弦関数という.\footnote{$\sinh x=\frac{e^x-e^{-x}}{2}$と合わせて,\,$\cosh^2x-\sinh^2x=1$を満たす.\,$\cos^2x+\sin^2x=1$の点$(\cos x,\sin x)$が円を描くのに対し,点$(\cosh x,\sinh x)$は双曲線$X^2-Y^2=1$を描くので,この名前がある.}
$u=e^x$と置けば$u+\frac1u=2\cosh x$であり,積が1になる2数$u,\frac1u$は$e^x,e^{-x}$として自然に現れる.
しかも,定義から直接計算すると
\[
  \cosh2x=2\cosh^2x-1
\]
が成り立つ.\,$\cos2x=2\cos^2x-1$とそっくりではないか.
両辺を2倍すれば$2\cosh2x=(2\cosh x)^2-2$となり,漸化式$a_{n+1}=a_n^2-2$そのものである.
そこで$a_n=2\cosh t_n$と置けば,\,$t_{n+1}=2t_n$という等比型になるのだ.
$2\cosh t_1=3$より$e^{t_1}=\frac{3+\sqrt5}{2}$なので,
\[
  a_n=2\cosh\left(2^{n-1}\log\frac{3+\sqrt5}{2}\right)
\]
を得る.\,$\cosh$を定義に戻せば,先ほどの一般項と一致する.
つまり基本対称式型は,三角関数型の倍角の公式を,\,$\cos$から$\cosh$に取り替えたものだったのである.

実数の範囲で整理しておこう.
\begin{itemize}
  \item $-2\le a_1\le2$なら$a_1=2\cos\theta$と書け,\,$a_n=2\cos(2^{n-1}\theta)$となる.(\sectionref*{節：三角関数型})
  \item $a_1\ge2$なら$a_1=2\cosh t$と書け,\,$a_n=2\cosh(2^{n-1}t)$となる.
  \item $a_1<-2$なら$a_1=-2\cosh t$と書ける.\,$a_2=a_1^2-2=2\cosh2t$で符号が消えるので,\,$n\ge2$では$a_n=2\cosh(2^{n-1}t)$である.
\end{itemize}

\paragraphbf{3倍角とチェビシェフ多項式}
同じことは3倍角でも起こる.
\begin{align*}
  &\left(u+\frac1u\right)^3\\
  &\qquad-3\left(u+\frac1u\right)=u^3+\frac{1}{u^3}
\end{align*}
なので,\,$a_1=3,\ a_{n+1}=a_n^3-3a_n$なら$a_n=2\cosh(3^{n-1}t_1)$である.
これは
\[
  \cosh3x=4\cosh^3x-3\cosh x
\]
に対応している.
一般に,\sectionref*{節：角度の整数倍と多項式}の\bookterm{chebyshev}{チェビシェフ多項式}$T_m$を使えば,
\[
  2T_m\left(\frac{u+u^{-1}}{2}\right)=u^m+u^{-m}
\]
が成り立つ.\footnote{$u=\cos\theta+i\sin\theta$のときは$T_m(\cos\theta)=\cos m\theta$そのものである.両辺は$u$の式として等しいので,すべての$u\ne0$で成り立つ.}
$a_{n+1}=a_n^2-2$は$m=2$,\,$a_{n+1}=a_n^3-3a_n$は$m=3$の場合であり,どちらも「$m$倍角の公式」として1つにまとめられるのだ.
$a_1=-3$のように初項が負のときは,\,$T_m$が偶関数か奇関数かによって符号の扱いが変わることに注意しよう.

\paragraphbf{行列のトレースで見る}
もう1つ,行列からの見方を紹介しよう.
正方行列の対角成分の和をトレースといい,\,$\operatorname{tr}$で表す.
$M=\begin{pmatrix}u&0\\0&u^{-1}\end{pmatrix}$とすれば$M^k=\begin{pmatrix}u^k&0\\0&u^{-k}\end{pmatrix}$なので,
\[
  u^k+u^{-k}=\operatorname{tr}(M^k)
\]
である.実は,\,$2\times2$行列$A$では
\[
  \operatorname{tr}(A^2)=(\operatorname{tr}A)^2-2\det A
\]
が成り立つ.\footnote{$A=\begin{pmatrix}a&b\\c&d\end{pmatrix}$で直接計算すれば,\,$\operatorname{tr}(A^2)=a^2+2bc+d^2$, $(\operatorname{tr}A)^2-2\det A=(a+d)^2-2(ad-bc)$となり一致する.}
$\det A=1$なら$\operatorname{tr}(A^2)=(\operatorname{tr}A)^2-2$であり,漸化式$a_{n+1}=a_n^2-2$は「行列式が1の行列を2乗したときのトレース」を追っていたのである.

\paragraphbf{小話：リュカ数}
$\frac{3+\sqrt5}{2}=\left(\frac{1+\sqrt5}{2}\right)^2$なので,\,$a_1=3$の場合の各項$3,7,47,2207,\ldots$は,黄金比$\varphi=\frac{1+\sqrt5}{2}$を使って$a_n=\varphi^{2^n}+\varphi^{-2^n}$と書ける.
$L_m=\varphi^m+\left(-\frac1\varphi\right)^m$で定まる数列はリュカ数と呼ばれ,フィボナッチ数列と同じ漸化式$L_{m+2}=L_{m+1}+L_m$を満たす.
$a_n$はリュカ数の$2^n$番目を順に取り出していたのだ.

\subsection{一次分数型と行列}\label{節：一次分数型と行列}

\paragraphbf{高校範囲の解法}
\[
  a_1=2,\qquad a_{n+1}=\frac{2a_n+1}{a_n+2}
\]
を\sectionref*{節：一次分数型}の方法で解こう.
特性方程式$x=\frac{2x+1}{x+2}$,すなわち$x^2=1$の2解は$\pm1$である.
$b_n=\frac{a_n-1}{a_n+1}$と置くと,公比は$\frac{2-1}{2+1}=\frac13$なので
\[
  b_{n+1}=\frac13b_n,\qquad b_1=\frac13,\qquad b_n=3^{-n}
\]
となる.\,$a_n=\frac{1+b_n}{1-b_n}$に戻して
\[
  a_n=\frac{1+3^{-n}}{1-3^{-n}}=\frac{3^n+1}{3^n-1}
\]
を得る.\footnote{$a_1=2>0$で,正の$a_n$から次の項も正になるので,分母$a_n+2$は0にならない.また$ps-qr=2\cdot2-1\cdot1=3\ne0$である.}

\paragraphbf{2つの量の線形な変化と,その比}
\sectionref*{節：分数を二つの成分で追う}で見たように,\,$a_n=\frac{x_n}{y_n}$と分子と分母に分けると,
\[
  \begin{pmatrix}x_{n+1}\\y_{n+1}\end{pmatrix}
  =\begin{pmatrix}2&1\\1&2\end{pmatrix}
  \begin{pmatrix}x_n\\y_n\end{pmatrix}
\]
となる.これは連立型
\[
  x_{n+1}=2x_n+y_n,\qquad y_{n+1}=x_n+2y_n
\]
そのものであり,\sectionref*{節：行列の基本}の最初に扱った例と同じ行列である.
$(x_1,y_1)=(2,1)$から和と差で解けば$x_n=\frac{3^n+1}{2},\ y_n=\frac{3^n-1}{2}$となり,その比が先ほどの$a_n$に一致する.
つまり,\textbf{一次分数型は,連立型で動く2つの量の比を追っていた}のである.
不動点$\pm1$は,行列の固有ベクトル$\begin{pmatrix}1\\1\end{pmatrix},\begin{pmatrix}-1\\1\end{pmatrix}$の成分比であり,公比$\frac13$は2つの固有値$1,3$の比である.
このように,ベクトルを行列で動かし,その成分の比だけを見る考え方を射影変換という.

この見方に立てば,連立型・隣接3項型・一次分数型の3つは,どれも2行2列の行列を何回も掛ける問題として1つにまとまる(\sectionref*{節：漸化式と行列}).
ただし,比を取る以上,分母にあたる成分$y_n$が0になる場合は別に扱う必要がある.

\subsection{総和を含む漸化式と母関数}\label{節：総和を含む漸化式と母関数}

\paragraphbf{高校範囲で分かること}
\[
  a_0=1,\qquad a_{n+1}=\sum_{k=0}^{n}a_ka_{n-k}
\]
を考えよう.まずは実験である.
\begin{align*}
  a_1&=1\cdot1=1,\\
  a_2&=1\cdot1+1\cdot1=2,\\
  a_3&=2+1+2=5,\\
  a_4&=5+2+2+5=14
\end{align*}
となり,\,$1,1,2,5,14,\ldots$と続く.
右辺の和は,添字の和が$n$で一定になる積を全部足したものである.
これは総和型のように$S_{n+1}-S_n$を取っても消えてくれず,高校範囲の変形では一般項にたどり着けない.

\paragraphbf{母関数で見る}
そこで\sectionref*{節：数列をまとめて見る}の母関数$A(x)=\sum_{n=0}^{\infty}a_nx^n$を使う.
2つの級数を掛けると,\,$x^n$の係数には添字の和が$n$になる積の和が現れる.
\[
  A(x)^2=\sum_{n=0}^{\infty}\left(\sum_{k=0}^{n}a_ka_{n-k}\right)x^n
\]
これはまさに漸化式の右辺である.したがって
\[
  A(x)=1+xA(x)^2
\]
という二次方程式が得られ,定数項が1になる解を選べば
\[
  A(x)=\frac{1-\sqrt{1-4x}}{2x}
\]
となる.ここまでは\sectionref*{節：数列をまとめて見る}で見た通りである.
「添字の和が一定となる積の総和」が,母関数の世界ではただの掛け算になるのだ.

\paragraphbf{係数を取り出す}
残るのは,この式から$x^n$の係数を取り出すことである.
平方根を展開するのは高校範囲を超えるので,ここでは別の道を通ろう.
$H(x)=\sqrt{1-4x}=\sum_{n=0}^{\infty}h_nx^n$と置く.\,$h_0=1,\ h_1=-2$である.
$H(x)$を微分すると$H'(x)=\frac{-2}{\sqrt{1-4x}}$なので,
\[
  (1-4x)H'(x)=-2H(x)
\]
が成り立つ.\footnote{ここでは級数を多項式と同じように項ごとに微分している.収束を考えず,係数の並びとして扱う形式的べき級数の立場では,この計算は正当化できる.}
両辺の$x^n$の係数を比べると,
\begin{align*}
  &(n+1)h_{n+1}-4nh_n=-2h_n\\
  &\therefore\quad h_{n+1}=\frac{4n-2}{n+1}h_n
\end{align*}
となる.なんと階比型である.
$A(x)=\frac{1-H(x)}{2x}$より$a_n=-\frac{h_{n+1}}{2}$なので,\,$a_n$も
\[
  a_n=\frac{4n-2}{n+1}a_{n-1}=\frac{2(2n-1)}{n+1}a_{n-1}
\]
という階比型を満たす.\sectionref*{節：階比型}と同じく書き下していけば,
\begin{align*}
  a_n&=\prod_{k=1}^{n}\frac{2(2k-1)}{k+1}\\
  &=\frac{2^n\cdot1\cdot3\cdots(2n-1)}{(n+1)!}\\
  &=\frac{(2n)!}{n!\,(n+1)!}
\end{align*}
すなわち
\[
  a_n=\frac{1}{n+1}\binom{2n}{n}
\]
を得る.\footnote{最後の変形では$(2n)!=\{1\cdot3\cdots(2n-1)\}\cdot\{2\cdot4\cdots(2n)\}=1\cdot3\cdots(2n-1)\cdot2^nn!$を使った.}
これがカタラン数である.
非線形で解けそうになかった漸化式が,母関数を経由すると,本編の最初に学んだ階比型に帰着したのだ.
母関数は,扱える問題の範囲を大きく広げてくれる道具なのである.

\subsection{多項式と指数の付加項と差分作用素}\label{節：多項式と指数の付加項と差分作用素}

\paragraphbf{高校範囲の解法}
\[
  a_1=0,\qquad a_{n+1}-2a_n=n2^n
\]
を考える.付加項$n2^n$の指数部分が,\,$a_n$の係数2と同じであることに注目しよう.
\sectionref*{節：階差等比複合型}にならって指数部分を取り除こう.\,$a_n=2^{n-1}b_n$と置いて両辺を$2^n$で割れば,
\[
  b_{n+1}-b_n=n
\]
という階差型になる.\,$b_1=0$から$b_n=\frac{n(n-1)}{2}$であり,
\[
  a_n=2^{n-1}\cdot\frac{n(n-1)}{2}=2^{n-2}n(n-1)
\]
を得る.
ここで気になるのは,付加項は$n$の1次式なのに,\,$b_n$が2次式になったことである.
\sectionref*{節：階差等比複合型}で不動関数を探したときは,付加項と同じ次数の式を引けばよかった.なぜ今回は次数が1つ上がったのだろうか.

\paragraphbf{差分作用素で見る}
\sectionref*{節：和と差から微積分へ}で導入した差分$\Delta b_n=b_{n+1}-b_n$を,数列を数列に移す操作と見て,差分作用素と呼ぶ.
多項式に$\Delta$を施すと,次数が1つ下がる.例えば
\begin{align*}
  \Delta n&=1,\\
  \Delta n^2&=2n+1,\\
  \Delta n^3&=3n^2+3n+1
\end{align*}
である.微分で$x^k$が$kx^{k-1}$になり次数が1つ下がるのと同じである.
逆に,\,$\Delta Q(n)=P(n)$となる$Q$を探すことは,微分すると$P$になる関数,つまり原始関数を探すことに対応する.
だから$Q$の次数は$P$より1つ上がるのだ.

一般に,\,$p,q\ne0$, $P(n)$を多項式として
\[
  a_{n+1}-pa_n=P(n)q^n
\]
を考えよう.\,$Q(n)q^n$の形の不動関数を探すと,代入して両辺を$q^n$で割り,
\[
  qQ(n+1)-pQ(n)=P(n)
\]
を満たす多項式$Q$を求める問題になる.
\begin{itemize}
  \item $p\ne q$なら,左辺の最高次の係数は$(q-p)$倍されるだけなので,\,$Q$は$P$と同じ次数で見つかる.
  \item $p=q$なら,左辺は$q\{Q(n+1)-Q(n)\}=q\Delta Q(n)$となり,差分を取ることになる.したがって$Q$は$P$より次数が1つ高くなければならない.
\end{itemize}
先ほどの例は$p=q=2$の場合であった.
特性根$p$と付加項の指数$q$が一致すると,次数が1つ上がる.これが「共鳴」と呼ばれる現象である.

隣接3項型などの高い階数の漸化式でも同じことが起こる.
特性方程式が$q$を$m$重解に持つと,\,$Q$の次数は$P$より$m$だけ上がり,\,$n^m$の因子が必要になる.
付加項のない隣接3項型で重解$\alpha$のとき一般項が$(A+Bn)\alpha^{n-1}$となったのも,同じ仕組みで説明できる.\footnote{添字を1つ進める操作を$E$と書くと,\,$a_{n+2}+pa_{n+1}+qa_n$は$(E^2+pE+q)a_n$と書ける.特性方程式が$(x-\alpha)^2$なら$(E-\alpha)^2(\alpha^nQ(n))=\alpha^{n+2}\Delta^2Q(n)$となるので,\,$\Delta^2Q=0$,すなわち$Q$は1次以下の多項式であればよい.}

\subsection{ニュートン法と変数変換}\label{節：ニュートン法と変数変換}

\paragraphbf{高校範囲の解法}
\sectionref*{節：ニュートン法}では$\sqrt2$を求めた.一般の$c>0$で
\[
  a_1>0,\qquad a_{n+1}=\frac12\left(a_n+\frac{c}{a_n}\right)
\]
を考えよう.これは方程式$x^2-c=0$に\bookterm{newton}{ニュートン法}を使ったときの漸化式である.
不動点は$\pm\sqrt c$であり,
\begin{align*}
  a_{n+1}-\sqrt c&=\frac{(a_n-\sqrt c)^2}{2a_n},\\
  a_{n+1}+\sqrt c&=\frac{(a_n+\sqrt c)^2}{2a_n}
\end{align*}
の比を取れば,\,$b_n=\frac{a_n-\sqrt c}{a_n+\sqrt c}$について$b_{n+1}=b_n^2$となる.
したがって$b_n=b_1^{2^{n-1}}$であり,
\[
  a_n=\sqrt c\cdot\frac{1+b_1^{2^{n-1}}}{1-b_1^{2^{n-1}}}
\]
を得る.\,$a_1>0$より$|b_1|<1$なので$b_n\to0$,すなわち$a_n\to\sqrt c$である.
誤差$e_n=a_n-\sqrt c$で見れば$e_{n+1}=\frac{e_n^2}{2a_n}$であり,誤差が2乗に比例して小さくなる.
なお,\,$a_1=\sqrt c$なら$b_1=0$であり,すべての項が$\sqrt c$の定数数列になる.

\paragraphbf{双曲線関数で見る}
$a_1>\sqrt c$のときは,\,$0<b_1<1$なので$b_1=e^{-2s}\ (s>0)$と書ける.
すると
\[
  \frac{1+e^{-2s}}{1-e^{-2s}}=\frac{e^s+e^{-s}}{e^s-e^{-s}}=\frac{\cosh s}{\sinh s}
\]
であり,\,$b_n=e^{-2^ns}$より
\[
  a_n=\sqrt c\cdot\frac{\cosh(2^{n-1}s)}{\sinh(2^{n-1}s)}
\]
となる.ここでも,角度にあたる$s$が毎回2倍されているのだ.\footnote{$0<a_1<\sqrt c$のときは$b_1<0$となり,この形では書けない.しかし相加平均と相乗平均の関係から$a_2\ge\sqrt c$となるので,\,$n\ge2$では同じ形で表せる.}

\paragraphbf{力学系の共役}
ここまでの5つの題材を振り返ると,共通の仕組みが見えてくる.
\begin{itemize}
  \item 基本対称式型では,\,$a=u+\frac1u$と置くと,\,$a\mapsto a^2-2$が$u\mapsto u^2$になった.
  \item ニュートン法では,\,$b=\frac{a-\sqrt c}{a+\sqrt c}$と置くと,\,$a\mapsto\frac12\left(a+\frac ca\right)$が$b\mapsto b^2$になった.
  \item 一次分数型では,\,$b=\frac{a-\alpha}{a-\beta}$と置くと,一次分数の操作が$b\mapsto\lambda b$になった.
\end{itemize}
どれも,ある変数変換$h$によって,難しい操作$f$を簡単な操作$g$に置き換えている.式で書けば
\[
  h(f(a))=g(h(a))
\]
である.このとき$f$と$g$は共役であるという.\footnote{漸化式のように,同じ操作を繰り返したときの振る舞いを調べる分野を力学系という.}
共役なら,\,$f$を$n$回繰り返すことは,\,$h$で移して$g$を$n$回繰り返し,\,$h$で戻すことと同じになる.
$g$が「2乗する」「定数倍する」のような既知の操作なら,一般項が求まるのだ.
本書を通して繰り返してきた「既知の形をつくる」という発想は,大学数学では共役という言葉で表されるのである.

\paragraphbf{まとめ}
5つの題材で,大学数学の見方がそれぞれ何を説明したのかをまとめておこう.
\begin{itemize}
  \item \textbf{思いつく理由}：基本対称式型の置き換え$u+\frac1u$は$2\cosh x$から,ニュートン法の比は2つの不動点から自然に現れる.
  \item \textbf{共通の仕組み}：一次分数型・連立型・隣接3項型は行列の累乗として,基本対称式型・三角関数型・ニュートン法は2倍の操作との共役として,1つにまとまる.
  \item \textbf{扱える問題の広がり}：母関数によって総和を含む非線形の漸化式が解け,差分作用素によって付加項の次数が上がる理由が分かる.
\end{itemize}
